% ======================================================================
%  sec02.tex —— 第 2 节 茎
%  底本：原书 pp.14–18（印刷页 = PDF 页）
%  原书标题：2. Stalks
%  体例：\section + \subsection 与原文编号一致；定理类环境紧跟 \label{类型:编号}。
% ======================================================================

\section{茎}\label{sec:2}

% ---------------------------------------------------------------- p.14
\subsection{定义}\label{sec:2.1}

\begin{definition}\label{def:2.1.1}
设 $X$ 是拓扑空间，$F$ 是 $X$ 上的集合预层。对 $x\in X$，令 $\mathcal U(X,x)\subseteq\mathcal U(X)$ 为 $X$ 中含 $x$ 的开邻域构成的全子范畴。$F$ 在 $x$ 处的\emph{茎}是
\begin{equation}\tag{2.1}\label{eq:2.1}
F_x:=\colim_{U\in\mathcal U(X,x)^{\op}}F(U).
\end{equation}
对给定的截面 $s\in F(U)$，$s$ 在 $x$ 处的\emph{芽}（记作 $s_x$）是 $s$ 在典范映射 $F(U)\to F_x$ 下的像。
\end{definition}

注意 $\mathcal U(X,x)^{\op}$ 是与正向集\footnote{正向集是指偏序集 $(I,\le)$，其中任意有限子集在 $I$ 中都有上界。} $(\mathcal U(X,x),\subseteq)$ 关联的范畴。故上述余极限是正向余极限\footnote{有人用「正向极限」（direct limit）一词。我强烈反对这一术语。}。由此 $F_x$ 可显式写成商
\begin{equation}\tag{2.2}\label{eq:2.2}
\Bigl(\coprod_{U\in\mathcal U(X,x)}F(U)\Bigr)\Big/\sim,
\end{equation}
即把一切 $F(U)$（$U\in\mathcal U(X,x)$）的无交并商掉等价关系 $\sim$。这里两个截面 $s\in F(U)$ 与 $s'\in F(U')$ 等价，当且仅当存在 $V\subseteq U\cap U'$ 使得 $s|_V=s'|_V$。用这一描述，截面 $s\in F(U)$ 的芽 $s_x$ 就是它所属的等价类。

\begin{remark}\label{rem:2.1.2}
一般地，设 $\mathcal C$ 是容许正向余极限的范畴，$F$ 是 $\mathcal C$-值预层。可用同一公式 \eqref{eq:2.1} 定义 $F$ 在 $x$ 处的茎。注意这一构造对 $F$ 函子化。

特别地，对阿贝尔群预层 $F$，可定义其茎 $F_x$，它是阿贝尔群。容易看出底层集合 $F_x$ 由 \eqref{eq:2.2} 给出，群结构由公式
\[
s_x+s'_x=(s|_V+s'|_V)_x,\qquad s\in F(U),\quad s'\in F(U'),\quad V\subseteq U\cap U'
\]
给出。
\end{remark}

% ---------------------------------------------------------------- p.15
\begin{remark}\label{rem:2.1.3}
设 $\mathcal C$ 是使其中正向余极限与有限极限交换的范畴\footnote{对 $\Set$、$\Ab$ 以及更一般的紧生成范畴成立。局部小范畴 $\mathcal C$ 中的对象 $c$ 称为\emph{紧}的，若 $\Hom_{\mathcal C}(c,-)$ 保持小滤过余极限。称 $\mathcal C$ 是\emph{紧生成}的，若它容许小余极限，且其中每个对象都同构于紧对象的小滤过余极限。}。则函子 $F\mapsto F_x$ 与有限极限交换。
\end{remark}

\begin{yizhu}
开邻域的全体在包含关系下构成正向集，故式 \eqref{eq:2.1} 是滤过余极限；它可以只沿任一共尾的邻域基求值，这正是「茎只看任意小的邻域」的确切含义。滤过余极限在 $\Set$ 与 $\Ab$ 中与有限极限交换，故茎函子保持有限极限。
[来源：nLab \emph{filtered colimit}，\url{https://ncatlab.org/nlab/show/filtered+colimit}（Last revised 2026-04-05；2026-11-17 访问）。]
\end{yizhu}

\subsection{层与茎}\label{sec:2.2}

下述结果表明，层的截面由它的芽决定。

\begin{lemma}\label{lem:2.2.1}
设 $F$ 是拓扑空间 $X$ 上的集合层。则对任意开子集 $U\subseteq X$，映射
\begin{equation}\tag{2.3}\label{eq:2.3}
F(U)\to\prod_{x\in U}F_x,\qquad s\mapsto (s_x)_{x\in U}
\end{equation}
是单射。此外，一族元素 $s(x)\in F_x$（$x\in U$）位于该映射的像中，当且仅当它满足下述条件：
\begin{itemize}
  \item[$(***)$] 对任意 $x\in U$，存在 $U$ 中含 $x$ 的邻域 $V$ 与截面 $s_V\in F(V)$，使得对任意 $y\in V$，有 $s(y)=(s_V)_y$。
\end{itemize}
\end{lemma}

\begin{proof}
先证映射 \eqref{eq:2.3} 是单射。设 $s,s'\in F(U)$ 的芽处处相等。按定义，对任意 $x\in U$ 存在 $V\subseteq U$ 使得 $s|_V=s'|_V$。特别地，可取开覆盖 $U=\bigcup_{i\in I}U_i$ 使 $s|_{U_i}=s'|_{U_i}$。但层条件蕴含
\[
F(U)\to\prod_{i\in I}F(U_i)
\]
是单射，故 $s=s'$。

映射 \eqref{eq:2.3} 的像中元素显然满足条件 $(***)$。反之，设 $s(x)\in F_x$（$x\in U$）是一族满足 $(***)$ 的元素。按假设，可取开覆盖 $U=\bigcup_{i\in I}U_i$ 与截面 $s_i\in F(U_i)$，使得对任意 $x\in U_i$ 有
\begin{equation}\tag{2.4}\label{eq:2.4}
s(x)=(s_i)_x.
\end{equation}
特别地，$s_i|_{U_i\cap U_j}$ 与 $s_j|_{U_i\cap U_j}$ 的芽处处相等。把 \eqref{eq:2.3} 的单射性施用于 $U_i\cap U_j$，得
\[
s_i|_{U_i\cap U_j}=s_j|_{U_i\cap U_j}.
\]
于是由层条件，存在唯一的 $s\in F(U)$ 使得 $s|_{U_i}=s_i$。对任意 $x\in U$，取 $i\in I$ 使 $x\in U_i$，则
\[
s_x=(s_i)_x=s(x),
\]
其中第一个等号来自茎的定义，第二个等号即 \eqref{eq:2.4}。特别地，$s(x)\in F_x$（$x\in U$）是 $s$ 在映射 \eqref{eq:2.3} 下的像。
\end{proof}

\begin{yizhu}
Stacks 把引理 \ref{lem:2.2.1} 的单射性单独命名为 separated presheaf（分离预层）：预层 $F$ 称为分离的，若 $F(U)\to\prod_{x\in U}F_x$ 对一切开集 $U$ 是单射。于是层条件被拆成「分离」与「粘接」两半，验证时可分别检查。
[来源：Stacks Project, Tag 007A（Definition 6.11.2，分离预层），\url{https://stacks.math.columbia.edu/tag/007A}（2026-11-17 访问）。]
\end{yizhu}

\begin{remark}\label{rem:2.2.2}
对预层的类似断言一般不成立。即取 $U=X=\emptyset$，空积 $\prod_{x\in\emptyset}F_x$ 是单元集，而 $F(\emptyset)$ 可以是任意集合。
\end{remark}

\begin{corollary}\label{cor:2.2.3}
若 $\alpha,\beta\colon F\to F'$ 是集合层之间的态射，且对任意 $x\in X$ 有 $\alpha_x=\beta_x$，则 $\alpha=\beta$。
\end{corollary}

\begin{remark}\label{rem:2.2.4}
用范畴论的语言，上述推论说的是：函子 $(-)_x\colon\Shv(X,\Set)\to\Set$（$x\in X$）联合保守。
\end{remark}

% ---------------------------------------------------------------- p.16
\begin{proposition}\label{prop:2.2.5}
设 $\alpha\colon F\to F'$ 是拓扑空间上集合层之间的态射。则 $\alpha$ 是同构，当且仅当对任意 $x\in X$，$\alpha_x\colon F_x\to F'_x$ 是双射。
\end{proposition}

\begin{proof}
「仅当」显然。下证「当」。设 $\alpha_x$ 对任意 $x\in X$ 都是双射。注意有交换图表
\[
\begin{tikzcd}
F(U) \arrow[r] \arrow[d,"\alpha_U"'] & \prod_{x\in U}F_x \arrow[d,"(\alpha_x)_{x\in U}"] \\
F'(U) \arrow[r] & \prod_{x\in U}F'_x.
\end{tikzcd}
\]
由引理 \ref{lem:2.2.1}，水平映射是单射，故 $\alpha_U$ 也是单射。

剩下证明 $\alpha_U$ 是满射。设 $s'\in F'(U)$ 是截面，我们来构造截面 $s\in F(U)$ 使 $s$ 在 $\alpha_U$ 下的像为 $s'$。

对任意点 $x\in U$，因 $\alpha_x$ 是双射，可取开子集 $V\subseteq X$ 与截面 $t\in F(V)$，使得 $\alpha_x(t_x)=s'_x$。按定义 $\alpha_x(t_x)=\alpha_V(t)_x$。故 $\alpha_V(t)$ 与 $s'$ 在 $x$ 处的芽相等。按定义，存在 $x$ 在 $U\cap V$ 中的开邻域 $W$，使得 $\alpha_V(t)|_W=s'|_W$。又注意 $\alpha_V(t)|_W=\alpha_W(t|_W)$。

由此可见，可取开覆盖 $U=\bigcup_{i\in I}U_i$ 与截面 $s_i\in F(U_i)$，使得 $\alpha_{U_i}(s_i)=s'|_{U_i}$。特别地，
\[
\alpha_{U_i\cap U_j}(s_i|_{U_i\cap U_j})=\alpha_{U_i\cap U_j}(s_j|_{U_i\cap U_j})=s'|_{U_i\cap U_j}.
\]
因已证 $\alpha_{U_i\cap U_j}$ 是单射，得 $s_i|_{U_i\cap U_j}=s_j|_{U_i\cap U_j}$。于是由 $F$ 的层条件，存在唯一的截面 $s\in F(U)$ 使得 $s|_{U_i}=s_i$。再用 $F'$ 的层条件，容易看出 $\alpha_U(s)=s'$。
\end{proof}

上述结果表明，层之间的态射由其在茎上诱导的映射决定。然而，层本身并不由它的茎决定。

\begin{exercise}\label{exc:2.2.6}
设 $X$ 是连通拓扑空间，$E\to X$ 与 $E'\to X$ 是次数相同的两个覆盖空间。证明：对任意点 $x\in X$，$X$ 上的层 $\operatorname{Sect}_E$ 与 $\operatorname{Sect}_{E'}$ 的茎同构；但除非存在定义在 $X$ 上的同胚 $E\simeq E'$，二者不同构。
\end{exercise}

\begin{yizhu}
此习题的要害在于：茎只记住「局部」信息，覆盖空间的整体拓扑（即基本群的作用）并不记录在茎里；因此纤维等势而层不同构是常态。反过来，茎此处的确只是纤维 $E_x$，这与例 \ref{ex:1.1.4} 中截面的定义相符。
[译者按（上下文推断，无外部依据）：连通性的假设用于保证各点纤维等势，否则连茎的基数都未必相同。]
\end{yizhu}

\begin{remark}\label{rem:2.2.7}
设 $\mathcal C$ 是紧生成范畴。引理 \ref{lem:2.2.1} 与命题 \ref{prop:2.2.5} 可推广到 $\mathcal C$-值层。换言之：
\begin{itemize}
  \item 对任意 $\mathcal C$-值层 $F$，态射 $F(U)\to\prod_{x\in U}F_x$ 是单态射。
  \item $\mathcal C$-值层之间的态射 $\alpha\colon F\to F'$ 是同构，当且仅当对任意 $x\in X$，$\alpha_x\colon F_x\to F'_x$ 是同构。
\end{itemize}
这些断言可由 $\Set$ 的特殊情形，借助下面两条观察推出：
\begin{itemize}
  \item $\mathcal C$ 中的态射 $d\to d'$ 是单态射（相应地，同构），当且仅当对任意紧对象 $c\in\mathcal C$，映射 $\Hom_{\mathcal C}(c,d)\to\Hom_{\mathcal C}(c,d')$ 是单射（相应地，双射）。
  % -------------------------------------------------------------- p.17
  \item 对任意 $\mathcal C$-值层 $F$ 与任意紧对象 $c\in\mathcal C$，集合值层
  \[
  \smallcat U(X)^{\op}\xrightarrow{\ F\ }\mathcal C\xrightarrow{\ \Hom_{\mathcal C}(c,-)\ }\Set
  \]
  在 $x\in X$ 处的茎，典范同构于 $\Hom_{\mathcal C}(c,F_x)$。
\end{itemize}
细节留给好奇的读者。
\end{remark}

\subsection{摩天楼层}\label{sec:2.3}

\begin{definition}\label{def:2.3.1}
设 $X$ 是拓扑空间，$x\in X$ 是点。对任意集合 $A$，如下定义集合预层 $\delta_{x,A}$。
\begin{itemize}
  \item 对开子集 $U\subseteq X$：
  \begin{itemize}
    \item 若 $x\in U$，定义 $\delta_{x,A}(U):=A$；
    \item 若 $x\notin U$，定义 $\delta_{x,A}(U):=\{*\}$。
  \end{itemize}
  \item 对开子集 $U\subseteq V$：
  \begin{itemize}
    \item 若 $x\in U$（从而 $x\in V$），定义限制映射 $\delta_{x,A}(V)\to\delta_{x,A}(U)$ 为 $\id_A$；
    \item 若 $x\notin U$，定义限制映射为唯一的映射 $\delta_{x,A}(V)\to\delta_{x,A}(U)=\{*\}$。
  \end{itemize}
\end{itemize}
容易验证这确实定义了预层 $\delta_{x,A}$。称它为在 $x$ 处取值 $A$ 的\emph{摩天楼层}。
\end{definition}

\begin{yizhu}
「摩天楼层」（skyscraper sheaf）之名来自图像：整栋楼只立在 $x$ 这一点，$x$ 的一切邻域都取到整栋楼 $A$，其余开集只剩一个点。nLab 独立收录了 skyscraper sheaf 一词，其标准构造正是沿闭点包含 $\{x\}\to X$ 的正像；本书例 \ref{ex:3.2.3} 将说明 $\delta_{x,A}\simeq i_*A$。
[来源：nLab \emph{skyscraper sheaf}，\url{https://ncatlab.org/nlab/show/skyscraper+sheaf}（Last revised 2016-06-23；2026-11-17 访问）。]
\end{yizhu}

\begin{exercise}\label{exc:2.3.2}
预层 $\delta_{x,A}$ 确实是层。
\end{exercise}

\begin{lemma}\label{lem:2.3.3}
设 $X$ 是拓扑空间，$x\in X$ 是点，$A$ 是集合。则 $\delta_{x,A}$ 在点 $y\in X$ 处的茎典范双射于
\begin{itemize}
  \item 集合 $A$，若 $y$ 含于 $\overline{\{x\}}$，即 $\{x\}$ 在 $X$ 中的闭包；
  \item 单元集 $\{*\}$，否则。
\end{itemize}
\end{lemma}

\begin{proof}
若 $y\in\overline{\{x\}}$，则 $y$ 的任意开邻域都含 $x$。于是
\[
(\delta_{x,A})_y:=\colim_{U\in\mathcal U(X,y)^{\op}}\delta_{x,A}(U)\simeq\colim_{U\in\mathcal U(X,y)^{\op}}A
\]
是取值 $A$ 的常图表的正向余极限，故 $(\delta_{x,A})_y\simeq A$。

若 $y\notin\overline{\{x\}}$，则存在 $y$ 的开邻域 $V$ 使 $x\notin V$。注意 $\mathcal U(V,y)^{\op}\subseteq\mathcal U(X,y)^{\op}$ 是（余）终的。于是
\[
(\delta_{x,A})_y:=\colim_{U\in\mathcal U(X,y)^{\op}}\delta_{x,A}(U)\simeq\colim_{U\in\mathcal U(V,y)^{\op}}\delta_{x,A}(U)\simeq\colim_{U\in\mathcal U(V,y)^{\op}}\{*\}
\]
是取值 $\{*\}$ 的常图表的正向余极限，故 $(\delta_{x,A})_y\simeq\{*\}$。
\end{proof}

\begin{yizhu}
点 $y$ 落在 $\overline{\{x\}}$ 中，等价于 $x$ 在特殊化序下特殊化到 $y$，记作 $x\leadsto y$；这是拓扑空间上由开集定义的预序。于是引理 \ref{lem:2.3.3} 可读作：摩天楼层的茎在 $x$ 的特殊化下闭上取 $A$、其外取单元集。对 $T_1$ 空间（如 $\Spec R$ 的闭点）闭包只有 $\{x\}$ 本身。
[来源：nLab \emph{specialization order}，\url{https://ncatlab.org/nlab/show/specialization+order}（Last revised 2023-06-06；2026-11-17 访问）。]
\end{yizhu}

注意若 $A$ 带有阿贝尔群结构，摩天楼层 $\delta_{x,A}$ 可提升为阿贝尔群层。此时阿贝尔群 $(\delta_{x,A})_y$ 或为 $A$ 或为 $0$。

\begin{proposition}\label{prop:2.3.4}
设 $X$ 是拓扑空间，$x\in X$ 是点，$A$ 是集合。则对任意集合预层 $F$，复合
\begin{equation}\tag{2.5}\label{eq:2.5}
\Hom_{\PShv(X,\Set)}(F,\delta_{x,A})\xrightarrow{\ (-)_x\ }\Hom_{\Set}(F_x,(\delta_{x,A})_x)\simeq\Hom_{\Set}(F_x,A)
\end{equation}
是双射。
\end{proposition}

\begin{corollary}\label{cor:2.3.5}
茎函子
\[
\PShv(X,\Set)\to\Set,\qquad F\mapsto F_x
\]
有右伴随
\[
\Set\to\PShv(X,\Set),\qquad A\mapsto\delta_{A,x}.
\]
\end{corollary}

% ---------------------------------------------------------------- p.18
\begin{proof}[命题 \ref{prop:2.3.4} 的证明]
先构造映射
\begin{equation}\tag{2.6}\label{eq:2.6}
\Hom_{\Set}(F_x,A)\to\Hom_{\PShv(X,\Set)}(F,\delta_{x,A})
\end{equation}
如下。给定映射 $f\colon F_x\to A$，对任意开子集 $U\subseteq X$ 定义映射 $\alpha_U\colon F(U)\to\delta_{x,A}(U)$：
\begin{itemize}
  \item 若 $x\in U$，令 $\alpha_U$ 为复合 $F(U)\to F_x\xrightarrow{\ f\ }A$；
  \item 若 $x\notin U$，令 $\alpha_U$ 为唯一的映射 $F(U)\to\{*\}$。
\end{itemize}
容易验证这些映射与限制映射相容，因此定义态射 $\alpha\colon F\to\delta_{x,A}$。现令映射 \eqref{eq:2.6} 为 $f\mapsto\alpha$。

容易验证 \eqref{eq:2.5} 与 \eqref{eq:2.6} 互为逆映射。故二者都是双射。
\end{proof}

\begin{remark}\label{rem:2.3.6}
一般地，对任意有终对象\footnote{对象 $*\in\mathcal C$ 是终对象，当且仅当对任意 $c\in\mathcal C$ 存在唯一的态射 $c\to *$。}的范畴 $\mathcal C$ 与任意对象 $A\in\mathcal C$，都可定义 $\mathcal C$-值层 $\delta_{x,A}$。若 $\mathcal C$ 容许正向余极限，则 $\delta_{x,A}$ 的茎或为 $A$、或为 $\mathcal C$ 的终对象，且函子 $A\mapsto\delta_{A,x}$ 右伴随于 $F\mapsto F_x$。
\end{remark}
